% 第 4 章：原生原对偶内点法（Parity IPM）
\section{原生原对偶内点法（Parity IPM）}\label{sec:ipm}

\subsection{功能定位与入口}
\compmeta{opf::parity::solve\_primal\_dual\_ipm}{\srcpath{include/hacdcpf/optimal\_power\_flow/native\_ipm\_solver.hpp}}
实现 src/optimal\_power\_flow/parity\_ipm.cpp:852--2149。对第~\ref{sec:parity}~章的
全空间 NLP 求 KKT 点：原变量 $x$、等式对偶 $\lambda$、不等式对偶 $\mu$、松弛 $z>0$，
障碍子问题
\[
\min\ f(x)-\gamma\sum_j\log z_j\quad\mathrm{s.t.}\quad g(x)=0,\quad h(x)+z=0,\quad x_{min}\le x\le x_{max}.
\]

\subsection{初始化与目标缩放}
不等式装配 \srcpath{assemble\_inequalities}（:799--844）：非线性行 + 有限上下界行。
松弛初始化 $z=\max(r_h,1)$ 或 $\max(-r_h,10^{-2})$，乘子 $\mu=\max(1/z,10^{-2})$
（:924--941）。目标缩放
\[
\fld{obj\_scale}=\min\bigl(1,\ 100/\lVert\nabla f\rVert_\infty\bigr)\quad(:957--987),
\]
治理大电网成本梯度淹没平衡行的问题（case2383wp 回归，见验证节）。

\subsection{收敛度量（归一化）}
:1014--1027：
\begin{align*}
\mathrm{feascond}&=\frac{\max(\lVert g\rVert_\infty,\ \max h^+)}{1+\max(\lVert x\rVert_\infty,\lVert z\rVert_\infty)},
&\mathrm{gradcond}&=\frac{\lVert\nabla f+J_g^{\mathsf T}\lambda+J_h^{\mathsf T}\mu\rVert_\infty}{1+\max(\lVert\mu\rVert,\lVert\lambda\rVert)},\\
\mathrm{compcond}&=\frac{z^{\mathsf T}\mu/n_{iq}}{1+\lVert x\rVert_\infty},
&\fld{comp\_tol}&=\max(10\,\mathrm{tol},\ 2\times10^{-3})\ (:1036).
\end{align*}
另保留未归一化的 \fld{raw\_feascond}。

\subsection{Newton 系统两种形式}
\srcpath{KktForm}（:167--199），env \texttt{HACDCPF\_OPF\_KKT\_FORM}，默认 condensed：
\begin{itemize}
  \item \textbf{Condensed}：$W=L_{xx}+J_h^{\mathsf T}\Sigma J_h$，$\Sigma=\mathrm{diag}(\mu/z)$
        （:1176--1189）；KKT
        $\begin{bmatrix}W+\delta I & J_g^{\mathsf T}\\ J_g & -\delta I\end{bmatrix}$（:443--481）；
        递进正则化 $\{0,10^{-8},10^{-6},10^{-4}\}\times\lVert W\rVert$（:1467--1473）。
        RHS 装配（:1461--1544）：预估 $N_{aff}=L_x+J_h^{\mathsf T}(\mu\circ r_h/z)$；
        校正 $N_{corr}=L_x+J_h^{\mathsf T}\big((\mu\circ r_h+\gamma e-\mathrm{cross})/z\big)$；
        恢复 $dz=-(r_h+z)-J_h\,dx$，
        $d\mu=-\big(\mu+\Sigma dz+(\mathrm{cross}-\gamma e)/z\big)$；
  \item \textbf{Augmented}：
        $\begin{bmatrix}L_{xx}+\delta_W I & J_g^{\mathsf T} & J_h^{\mathsf T}\\
        J_g & -\delta_C I & 0\\ J_h & 0 & -ZM^{-1}\end{bmatrix}$（:562--629）；
        Wächter--Biegler 惯性校正：MUMPS 负主元数须 $=n_{eq}+n_{iq}$，否则 $\delta_W$
        从 $10^{-8}\lVert L_{xx}\rVert$ 起步 $\times8$ 增至 $10^{-2}\lVert L_{xx}\rVert$ 上限
        （:1281--1301）；$\delta_C$ 默认 $10^{-8}\lVert L_{xx}\rVert$（:1232--1239，
        可由 \texttt{HACDCPF\_OPF\_DELTA\_C} 绝对值覆盖）。
        \textbf{$\delta_W$ 暖启动与 KLU/MUMPS 不对称}（:1271--1319）：
        仅 MUMPS 步用 $\delta_W=\max(10^{-8}\lVert L_{xx}\rVert,\ 0.5\,\delta_{W,prev})$
        （Ipopt $\kappa_W^-$ 规则）；KLU 步恒 $\delta_W=0$（惯性不可报，且防止
        $\delta_W$ 泄漏污染 KLU 轨迹，实测 case13659 +44 迭代）；
        惯性循环最多 8 次、第 8 次封顶接受；\texttt{HACDCPF\_OPF\_INERTIA\_GATE}=N
        门控的周期 MUMPS 复检后\textbf{总是恢复 KLU}。
\end{itemize}

\subsection{线性代数后端与数值健壮性}
后端优先级 \srcpath{backend\_order}（:268--327）：结构感知——纯 AC 先 KLU
（BTF 结构实测最快），混合 AC/DC 先 MUMPS（需 $LDL^{\mathsf T}$ 惯性）；
UMFPACK/Eigen SparseLU 为升级链；env \texttt{HACDCPF\_OPF\_LINEAR\_SOLVER} 可钉死；
Windows 上 KKT 维数 $\le1024$ 自动走 dense PartialPivLU（:1069--1078）。
健壮性措施：对称 Ruiz 均衡 5 次扫描、每次求解冻结（:415--438, 503--517）；
2 步迭代精化带发散守卫（:655--671）；Eigen SparseLU 残差超
$10^{-6}\lVert rhs\rVert$ 标记 \fld{solve\_degraded} 触发后端升级（:677--683）。

\subsection{步长计算与全局化}
\begin{itemize}
  \item \textbf{Mehrotra 预估-校正}：预估 $\gamma=0$，$\sigma=(\mu_{aff}/\mu)^3$
        （:1500--1507）；校正 cross 项限幅 $\pm0.1\gamma$（:1513--1518）；
        Gondzio 附加校正 $\le2$ 次（:1554--1591）；
        \textbf{校正器再中心化重试}（:1375--1450）：校正方向若 ftb 步长 $<10^{-9}$
        （$\kappa\sim1/\mu$ 爆炸征兆），则 $\gamma_{corr}=\max(\gamma,\ 0.5\,\mu_{cur})$、
        cross 清零重解，$\le2$ 次。
        \textbf{ABO 动态 $\mu$ 策略}（可选，env \texttt{HACDCPF\_OPF\_MU\_STRATEGY}，
        :1343--1359, 1454--1456）：$\mu^+=\mu^2/\mu_0$（$\theta_S$ 更新，$\gamma_q=1$，
        $b=\mu_0$），$\gamma=\mathrm{clamp}\big(\mu+\alpha_{aff}(\mu^+-\mu),\ 10^{-14},\ 10^4\big)$，
        $\alpha_{aff}=\min(\alpha_{p,aff},\alpha_{d,aff},1)$；
        $\mu_{dyn}/\mu_0$ 首迭代初始化为当前 $\mu_{cur}$，校正步后
        $\mu_{dyn}=\gamma_{corr}$（含再中心化）；
  \item \textbf{$(\theta,\varphi)$ 滤波线搜索}（Wächter--Biegler §3.2）：
        常数 $\gamma_\theta=\gamma_\varphi=10^{-5}$、$\Delta=1$、$s_\theta=1.1$、
        $\eta_\varphi=10^{-8}$、$\theta_{min}=10^{-4}$、
        $\theta_{max}=10^4\max(1,\text{初始不可行度})$（:1625--1641）；
        domination 记忆 \fld{filter\_history}（:1731--1776）；
        $\varphi_\gamma=\tilde f-\gamma\sum\log z$（:1669--1681）；
        仅当 KKT 维数 $\ge512$ 或 env \texttt{HACDCPF\_OPF\_FILTER\_ALL} 启用（:1646--1653）；
        另有 legacy 三轴接受路径（feas/stationarity/complementarity 进展，
        带对偶轴 $\times100$ 防爆守卫，:1707--1729）；
  \item \textbf{其他}：二阶校正 SOC（仅 augmented，精确触发：bt==0 且
        $\theta_{trial}>\theta_k$ 且 meq$>$0，rhs $=-(\alpha r_g+r_g(x+\alpha dx))$，
        :1800--1874）、中心性恢复重试（:1890--1907）、恢复阶段 $\rho I$
        Gauss--Newton 可行性步（见下节）、
        空步规则 $\min(\alpha_p,\alpha_d)<10^{-7}$ 判方向不可用（:1600--1608）。
\end{itemize}
另：\fld{obj\_scale} 的 Hessian 修正 $(\fld{obj\_scale}-1)\cdot h_{diag}$ 加在
$L_{xx}$ 对角（:973--987）。

\subsection{终止、认证与同伦切线}
\begin{itemize}
  \item \textbf{best-iterate 跟踪}（:1993--2005）与\textbf{可接受收敛}
        （:2038--2074）：feas/raw/grad $<100\times$tol 且 comp $<$ \fld{comp\_tol}，
        best-vs-final 择优，status 区分 ``best iterate''/``final iterate''；
  \item \textbf{目标停滞提前停止}（:2007--2026）：$\fld{raw\_feas}<10^{-4}\ \land\
        \mathrm{comp}<\fld{comp\_tol}$ 且 $|\Delta obj|\le10^{-6}(1+|obj|)$
        连续 20 接受步，status=``objective-stagnant (dual-polish certification pending)''；
  \item \textbf{恢复阶段}（:1909--1962）：$(1,1)$ 块 $\rho I$，
        $\rho=10^{-8}\lVert L_{xx}\rVert$，$\delta_C=10^{-8}\lVert L_{xx}\rVert$，
        尾部 rhs 用 $\gamma_{corr}$；$\le8$ 次 $\alpha$ 减半，接受条件 $\theta$ 改善
        $(1-10^{-4}\alpha)$ 且对偶轴 $\times100$ 有界；
  \item \textbf{对偶最小二乘抛光}（:2076--2126）：解
        $\begin{bmatrix}I & J_g^{\mathsf T} & J_h^{\mathsf T}\\ J_g & 0 & 0\\
        J_h & 0 & 0\end{bmatrix}(r,\lambda,\mu)=(-\tilde f\nabla f,\ 0,\ 0)$
        （$z=0$、$\mu=e$ 的增广 KKT，$\delta_C=10^{-8}$），$\mu$ 裁 $\ge0$；
        采用规则 $\fld{stat\_ls}\le\fld{gradcond}$ 或未收敛即采用；
        认证 status=``converged (dual least-squares certified)''；dense 路径跳过。
\end{itemize}
\srcpath{homotopy\_tangent}（:2154--2257）：一次惯性控制 KKT 求解
$K_{aug}\,(dx/dt,d\lambda/dt,d\mu/dt)=(-\nabla f/t,\ 0,\ 0)$，$dz/dt=-J_h\,dx/dt$，
供目标同伦 Davidenko 预测（第~\ref{sec:acopf}~章）。

\subsection{参数表（IPMOptions）}
定义 native\_ipm\_solver.hpp:17--38。
\begin{paramtable}
\fld{max\_iter} & -- & 次 & 100--1000 & 400 & IPM 最大迭代数（AC OPF 映射为 inner$\times$outer）\\
\fld{tol\_primal} & -- & -- & $10^{-8}$--$10^{-4}$ & $10^{-6}$ & 原始可行容差\\
\fld{tol\_dual} & -- & -- & $10^{-8}$--$10^{-4}$ & $10^{-6}$ & 对偶（stationarity）容差\\
\fld{tol\_complementarity} & -- & -- & $10^{-8}$--$10^{-4}$ & $10^{-6}$ & 互补容差（映射自 \fld{barrier\_mu\_min}）\\
\fld{regularization} & $\delta$ & -- & $10^{-10}$--$10^{-4}$ & $10^{-8}$ & KKT 递进正则化起步\\
\fld{alpha\_max} & $\alpha_{max}$ & -- & 0.9--0.999 & 0.95 & 最大步长（映射自 \fld{interior\_fraction}）\\
\fld{verbose} & -- & bool & -- & false & 迭代日志\\
\fld{primal\_start} 等四项 & -- & 向量 & -- & 空 & 暖启动 $x/\lambda/\mu/z$\\
\end{paramtable}
IPMResult 字段（converged、三项残差及初始值、factorization/linear\_solve/accepted/
rejected/backend\_escalations 计数、\fld{linear\_solver} 后端名等）见第~\ref{sec:options}~章。

\subsection{验证与校核}
\begin{itemize}
  \item case300\_acdc：accepted$>0$ 且 $\le$iterations、rejected$>0$
        （test\_opf\_solver\_backends.cpp:79--105）；
  \item case2383wp 目标缩放回归：修复前 stationarity$\approx0.45$ 顶格，现收敛（:481--494）；
  \item case2869pegase 稀疏 KKT $<100$ 迭代（:577--599）；case6515rte 保持有限不发散（:601--618）；
  \item dense KKT 后端（\texttt{HACDCPF\_OPF\_LINEAR\_SOLVER=dense}）解 case9/case30（:555--573）；
  \item 大系统诊断与 filter 回溯修复过程记录于 \srcpath{docs/large\_hybrid\_ipm\_diagnostics.md}
        （$\ge512$ 维 KKT 启用残差 filter 回溯，0.5 回溯最多 14 次）。
\end{itemize}
